
Large language models produce confident outputs regardless of factual accuracy. Output-level uncertainty metrics such as token entropy correlate weakly with correctness, achieving approximately 0.62 AUROC on factual question-answering benchmarks. Multi-sample methods such as semantic entropy improve detection to approximately 0.80 AUROC but require 5--20 forward passes, limiting practical deployment. This work investigates whether middle-layer hidden states encode a correctness signal that can be extracted via linear probing. Training a logistic regression probe on layer-15 (50\% depth) representations from Llama-3-8B-Instruct, we achieve 0.885 AUROC on TriviaQA, exceeding token entropy by +26.2 AUROC points with a single forward pass. Systematic layer sweeps across 8 depths (12.5\%--100\%) reveal an inverted-U pattern: correctness signal peaks at middle layers where semantic representations are richest, then declines in later layers where output formatting dominates. These results demonstrate that transformer hidden states encode factual accuracy information distinct from output confidence, enabling efficient correctness detection for LLM deployment.

